\documentclass[12pt]{article}
\usepackage[dutch]{babel}
\usepackage[onecolumn,portrait,a5paper,margin=14mm]{geometry}

\usepackage[twocolumn,landscape,a4paper,margin=15mm,columnsep=30mm]{geometry}
\usepackage{graphicx}
\usepackage{parskip}
\usepackage{color}
\usepackage{PTSansNarrow} 
\usepackage[T1]{fontenc}
\usepackage[dutch]{babel}
\usepackage{array}
\renewcommand\familydefault\sfdefault

\pagestyle{empty}

\renewcommand\emph[1]{{\slshape#1}}

\definecolor{DolRed}{rgb}{0.92,0.11,0.15} % #eb1c26
\definecolor{DolGray}{rgb}{0.82,0.82,0.82}

\newcommand\DolHeader[1]{
\includegraphics[width=.45\columnwidth]{dol_logo.png}\hfill\parbox[b]{.5\columnwidth}{#1}
\DolHline \\
{\bf Gaffelstraat 61B \DolSep Toegang gratis \DolSep Inloop: 19:30 uur \DolSep Start: 20:00 uur \DolSep Eind: 22:00 uur} \\
\DolHline
\vfill
}

\newcommand\DolHeaderEn[1]{
\includegraphics[width=.45\columnwidth]{dol_logo.png}\hfill\parbox[b]{.5\columnwidth}{#1}
\DolHline \\
{\bf Gaffelstraat 61B \DolSep Free entrance \DolSep Doors open: \ 19:30 \DolSep Start: \ 20:00 \DolSep End: \ 22:00} \\
\DolHline
\vfill
}

\newcommand\DolFooter{
\vfill
\colorbox{DolRed}{\parbox{\columnwidth}{\ \parbox[b]{.88\columnwidth}{\color{white}\raggedright
{\bf `Denken over links', een maandelijkse lezingenserie in het SP-pand aan
de Gaffelstraat. Doe mee, denk mee! Iedereen is welkom.} \\[-7pt]
\hfill denkenoverlinks.nl\\[.2em]}
\hfill \includegraphics[scale=.4]{qr.png}}}
}

\newenvironment{DolQuote}
{%
  \begin{quote}
  \begin{picture}(0,0)\put(-20,-34){\fontsize{2cm}{1em}\selectfont\bf\color{DolGray}``}\end{picture}
  \raggedright
  \setlength\parindent{1.5ex}
}{%
  \end{quote}
}

\newcommand\largeblack[2]{{\fontsize{#1}{1em}\selectfont\bf#2}\par}

\newcommand\DolHline{\textcolor{DolRed}{\rule{\columnwidth}{6pt}}}
\newcommand\DolBullet{\textcolor{DolGray}{$\bullet$}}
\newcommand\DolSep{\hfill\DolBullet\hfill}
\newcommand\DolMega[1]{{\fontsize{1.15cm}{1em}\selectfont\bf#1}\par}
\newcommand\DolHuge[1]{\textcolor{DolRed}{\Huge\bf#1}\par}
\newcommand\DolLarge[1]{\textcolor{DolRed}{\Large\bf#1}\par}
\newcommand\DolBox[1]{\colorbox{DolRed}{\parbox{\columnwidth}{\centering\parbox{.95\columnwidth}{\rule{0pt}{13pt}\color{white}\large\bf#1\rule[-7pt]{0pt}{0pt}}}}}


\begin{document}

\DolHeaderLarge{ \raggedleft
\largeblack{6.5mm}{
{\bf Donderdag 27 november} \\[7pt]
Rotterdam \\[7pt]}
}

\emph{Denken over Links is een lezingenserie waarin het denken over hedendaagse
sociale vraagstukken wordt geprikkeld en gescherpt. Deze maand:}

\vfill

\largeblack{10.6mm}{\bf STAD ZOEKT MIDDENKLASSE}

Met als sprekers o.a.

\textbf{Bart van Bouchaute}

\hfill\parbox{.95\columnwidth}{
Politicoloog verbonden aan het departement Sociologie van de Universiteit
Antwerpen, onderzoekscentrum Ongelijkheid, Armoede, Sociale Uitsluiting. Hij
is auteur van o.a. het boek `Gentrificatie als strategie van stadsvernieuwing'.}

\textbf{Joke van der Zwaard}

\hfill\parbox{.95\columnwidth}{
Joke van der Zwaard is ontwikkelingspsycholoog en werkt als zelfstandig
onderzoeker/publicist op het terrein van sociale ongelijkheid, sociale
mobiliteit, emancipatie en samenlevingskwesties in relatie tot wonen, opvoeding
en onderwijs. Ook is ze initiatiefnemer van o.a. De Leeszaal Rotterdam West.}

De sprekers gaan in op de voor- en nadelen van het trekken van `kansrijken'
naar de stad. Versterken zij de `sociale mix' in de stad en wijken?  En wat
betekent deze verandering voor stad en stedelingen?

\DolFooterLarge

\end{document}
